\textbf{T3: Plate between magnets (10 pts)}

Two identical long, cylindrical rod magnets of radius $R$ are close to each
other and share the same vertical symmetry axis. The polarity of the two
magnets is the same. As a result, the magnetic field in the air gap between
the magnets is directed towards the $+z$ direction (see \textit{figure}) and
uniform with flux density $B$. The magnetic field outside the gap is zero. A
horizontal large non-magnetic metal plate is placed in the air gap and moved
with constant horizontal velocity $v$ in $+y$-direction. The thickness of the
plate is $\delta$, the resistivity of the metal is $\varrho$.

\begin{center}\includegraphics[width=0.44\linewidth]{figures/EuPhO_2023_Q3_fig1.png}\end{center}

a) (3 pts) Sketch the shape of current streamlines in the metal plate at a
given time. Indicate the axes on your sketch.

b) (5 pts) Find and plot the current density inside the plate along a line
parallel to the $y$ axis intersecting the symmetry axis of the magnets.

c) (2 pts) Find the horizontal force required to move the plate.
